%% notes-tex/030-scaling-plots/sections/2-surface.tex

\section{The loss surface, read both ways\secstatus{todo}}
\label{sec:surface}

The grid is a five-by-four sample of $L(N, D)$ at one seed, with the two end
sizes replicated at the two smallest subsets. That is what makes the surface
family possible: every row of Figure~\ref{fig:surface} is a cut of the same
surface, and the two single-axis fits of the study are its edge columns.

\tcfig{figures/030-scaling-plots-surface.pdf}{%
The surface family, sketched. \textbf{a}, S1: loss against $D$ at every size,
one curve per size. Parallel curves on log axes are one data exponent for the
family; a fan is a data exponent that depends on size, which the single-column
fit of the study would not see. The replicate floor is dashed. \textbf{b}, S2:
loss against $N$ at every subset, one curve per subset. Each curve flattens at
the plateau its data sets, $E + B D^{-\beta}$, dashed per curve, and the size at
which a curve meets its plateau is the size beyond which more parameters stop
paying at that $D$. \textbf{c}, S3: the same twenty cells as a heatmap of the
matched read, with a second heatmap for the selected read in the real figure.
\textbf{d}, S4: the fitted forms drawn on $L - E$, where the floored form is a
straight line, with the bootstrap band and the held-out point ringed. Every
curve is hand-drawn to show the shape the plot would hold.
}{fig:surface}
\src{notes/assets/drawio/030-scaling-plots.gen.py}

\notesec{S1, loss against data at every size}
Five curves, each a row of the grid, on both reads and per screen. The study fits
the 180x8 curve alone; the other four are the control on that fit. If the small
sizes' curves run parallel to it, the data exponent is a property of the data
and the single fit stands for the family. If the small curves flatten earlier,
the small sizes are capacity-limited at large $D$, and the exponent fit at 180x8
is the largest the family shows, not the family's. The 1/64 and 1/16 points of
the end sizes carry three seeds each, so the spread of the replicated cells is
drawn as the noise floor of every curve.

\notesec{S2, loss against size at every subset}
Four curves, each a column. On the matched read the curves are directly
comparable, since every cell is read at the same records seen. On the selected
read the full-pool curve is a ten-epoch slice and sits above where convergence
would put it, so the plateau it appears to meet is partly the budget, which is
the limit that keeps this plot partial. The plateau lines are the second term of
the floored form fitted on each column, and the gap between a plateau and the
replicate floor is what more data could still buy at that size.

\notesec{S3, the heatmap}
Twenty cells shaded by the matched read, with the selected read beside it. The
reading is local: at any cell the steeper descent, down the column toward
larger $N$ or along the row toward larger $D$, is the direction the next unit of
effort should take from that cell. The two heatmaps differ most at the full-pool
column, and the size of that difference is a picture of how far the ten-epoch
budget sits from convergence.

\notesec{S4, the fitted forms}
Each axis of the study drawn on $L - E$ for the floored form, so the fit is a
straight line, with the pure power law beside it and the 5 to 95 percent band
from 300 resamples over runs. The held-out largest point is ringed with the
prediction the in-sample fit made for it. On the $D$ axis there are four points,
so the floored held-out test is undefined (three points fit three parameters)
and only the power form is tested there; that is the partial status.
